\section{课程知识体系总览}

讲者在课程结尾回顾了整个课程到目前为止的知识发展脉络。

\subsection{技术演进路径}

从课程开始到第 10 讲，讲者构建了以下完整的技术脉络：

\begin{enumerate}
    \item \textbf{VAE/GAN}（Lecture 1--2）：单步生成模型，质量和多样性受限
    \item \textbf{DDPM}（Lecture 3--4）：引入多步迭代去噪，大幅提升生成质量
    \item \textbf{DDIM}（Lecture 5--6）：DDPM 的确定性采样变体，支持跳步
    \item \textbf{SDE/ODE}（Lecture 7）：连续时间框架，统一视角
    \item \textbf{DPM-Solver}（Lecture 8）：利用 ODE 结构加速采样
    \item \textbf{Consistency Model}（Lecture 9）：学习快捷映射实现少步生成
    \item \textbf{Flow Matching}（Lecture 9--10）：更一般的框架，Rectified Flow + Reflow 实现直线轨迹和单步生成
\end{enumerate}

\begin{warningbox}{Flow Matching 不是扩散模型的``替代品''}
虽然本讲强调了流模型的优势，但需要注意：
\begin{itemize}
    \item 扩散模型和流模型在数学上是等价的（在线性高斯条件路径下）
    \item 已训练的扩散模型可以被直接转换为流模型使用
    \item 选择哪种框架更多取决于实际考量：训练稳定性、采样效率、代码复杂度等
    \item 当前的``最佳实践''是使用 Flow Matching 框架 + Rectified Flow 训练 + Reflow 微调
\end{itemize}
\end{warningbox}

\subsection{后续课程预告}

讲者预告了接下来三周的课程内容：
\begin{itemize}
    \item \textbf{下一周}：Inference-time Guidance（推理时引导），如何在不重新训练的前提下注入用户偏好
    \item \textbf{之后}：Discrete Diffusion Models（离散数据的扩散模型），将扩散/流的思想拓展到非连续数据（如文本）
\end{itemize}

\subsection{本章小结}

从 GAN 到扩散模型再到流模型，生成模型领域经历了从单步到多步再到少步的演进。Flow Matching 提供了一个统一且更灵活的框架，既能理解扩散模型的已有成果，又能为未来的改进提供空间。

